% Research progress snapshot, 7 November 2026; numerical sources remain in source comments.
\addchap{PART 1. INTRODUCTION}\label{chap:introduction}
\ThesisFlatNumbering

\section{Problem Statement and Scientific Motivation}
\subsection{From Proteins to Regulated Degradation}
Proteins carry out catalysis, transport, signalling and structural functions in cells. Their usefulness
depends on having the appropriate amount of functional protein in the appropriate place and at the
appropriate time. \emph{Protein homeostasis}, or proteostasis, coordinates protein production, folding,
localisation and removal; an imbalance can leave insufficient functional protein or allow damaged or
unneeded protein to accumulate \citep{balch2008}. Degradation therefore contributes both to quality
control and to regulation. The lifetime of a protein is one component of its abundance, alongside its
rate of synthesis; abundance and degradation rate are related quantities with different meanings.

The ubiquitin--proteasome system (UPS) provides one route for selective protein removal. A
\emph{degron} is a degradation signal whose activity depends on recognition and molecular context.
Such signals can involve N-terminal, C-terminal or internal sequence features; they are not a single
universal motif \citep{varshavsky2019,koren2018}. Changes to an amino acid can alter a recognition
motif or the accessibility of a signal. An exposed peptide sequence in a reporter and the same sequence
buried in a folded protein can consequently behave differently. Proteome-scale peptide measurements
indeed identify potent degrons in regions that are buried in the native structure \citep{voutsinos2026}.

\subsection{Mutation Measurements and Their Limits}
Experimental mutagenesis links sequence changes to measured phenotypes. Scanning mutagenesis tests
selected replacements along a sequence, whereas saturation mutagenesis samples a broader set of
replacements at each position. Early massively parallel sequence--function experiments showed that a
few representative replacements can miss important amino-acid-specific behaviour \citep{fowler2010}.
GPS profiling measures fluorescent fusion reporters in cells \citep{yen2008}; a later application to
internal degrons combined peptide screening, mutagenesis and mechanistic follow-up
\citep{zhang2023}. These approaches make mutation effects observable, but their output is tied to the
assay, reporter and cellular condition used.

The practical need is supported by experimental studies rather than by an assumed shortage of
computing. For example, VAMP-seq quantified the abundance of 7,801 single-amino-acid variants of
PTEN and TPMT, while also documenting proteins for which the assay was unsuitable
\citep{matreyek2018}. Zhang et al.\ report that a short-peptide assay can miss conformational motifs,
some modification-dependent signals and signals whose recognising factors are unavailable in the
assayed cells \citep{zhang2023}. Thus, extensive measurements still leave gaps in sequence coverage
and biological context. A computational model could help rank additional substitutions for follow-up.
The equal-budget OOF replay in Section~\ref{sec:res-prioritization} quantifies how well stored
scores recover measured hits. Actual laboratory savings require prospective testing: CT553H has
measured neither a reduction in wet-lab cost nor an improvement in prospective experimental yield.

\subsection{The Specific Prediction Problem}
This project predicts the change in the Protein Stability Index (PSI) of a peptide reporter after one
amino acid substitution:
\begin{equation}
 f(s,i,a)\longrightarrow\widehat{\Delta\mathrm{PSI}},\qquad
 \Delta\mathrm{PSI}=\mathrm{PSI}_{\mathrm{mutant}}-\mathrm{PSI}_{\mathrm{WT}},
 \label{eq:intro-task}
\end{equation}
where $s$ is the wild-type (WT) peptide sequence, $i$ is the substituted position and $a$ is the new
amino acid. In the Zhang assay, higher PSI means greater reporter stability. A positive effect therefore
indicates reporter stabilisation and may be consistent with weakening a degradation signal; it does
not identify the responsible E3 ligase or establish the mechanism. The task is continuous
mutation-effect regression, with event ranking as a secondary analysis. It is distinct from identifying
degron-containing peptides, predicting an absolute activity score, estimating folding free energy or
measuring the half-life of a natural full-length protein.

\section{Related Work and the Research Question}
\subsection{What Existing Studies Establish}
Experimental maps cover internal degrons \citep{zhang2023}, cytosolic peptides
\citep{voutsinos2026} and transcription-factor peptides \citep{larsen2025}; the last source is cited here
as the inspected 2025 preprint. PAP predicts peptide abundance with a convolutional model and uses
mutant--WT score differences to investigate variant abundance effects, including selected experimental
validation \citep{voutsinos2026}. MetaDegron predicts E3-targeted degrons using sequence and other
features \citep{zheng2024}. DEGRONOPEDIA and DegronMD provide degron prediction and contextual
annotation resources \citep{degronopedia,degronmd}. Degformer is an inspected software implementation
of transformer-based peptide prediction, cited as a repository rather than as a verified peer-reviewed
article \citep{degformer}. These works demonstrate that computational degron prediction and analysis
of mutation effects already exist.

Consequently, combining a protein language model (pLM) with a CNN is not, by itself, an established
novel contribution. The narrower question in CT553H is how well continuous substitution effects can be
predicted on a specified reporter endpoint when training-label coverage, sequence representation,
head architecture and adaptation are compared under common protein-group-aware evaluation.
Biological dependencies can otherwise make prediction estimates optimistic
\citep{bernett2024}. This is a bounded empirical question on the available data, not a claim that prior
work has never studied residue effects or controlled sequence similarity.

\subsection{Research Questions Organising This Progress Report}
The following questions connect the existing protocols and analyses. They organise the present
report; they are not presented as an independent preregistration preceding all historical experiments.
\begin{enumerate}
  \item \textbf{RQ1---Training-label coverage:} Does expanding the WT-PSI range used for training improve
        prediction on a common evaluation population, and which PSI strata benefit or deteriorate?
  \item \textbf{RQ2---Representation and head:} Do ordered context and pLM representations improve
        matched baselines, and does CNN or multi-scale CNN mixing add a reproducible advantage over
        a pointwise control?
  \item \textbf{RQ3---Transfer learning:} Under which tested conditions does scanning-mutation or
        peptide-label pre-training improve saturation mutation-effect prediction?
  \item \textbf{RQ4---Adaptation:} Does LoRA improve over its matched frozen-encoder control across
        seeds and PSI strata, and what computational cost accompanies the gain? Full fine-tuning
        requires its own completed evidence and is not inferred from LoRA.
  \item \textbf{RQ5---Generalisation and reliability:} How stable are predictions across held-out protein
        groups, where do large errors occur, and how do regression and probability calibration behave
        on a held-source population?
\end{enumerate}
The remaining literature search on work in Vietnam is \emph{not yet completed}; no claim of national
novelty is made. Independent external and biological validation remain open research needs.

\section{Objectives}
\subsection{General Objective}
To evaluate methods for predicting reporter-based amino acid substitution effects in degron peptides,
with explicit control of sequence relatedness and uncertainty, and to support traceable exploration
through a local software tool.

\subsection{Specific Objectives}
\begin{enumerate}
  \item Build an auditable pipeline for source provenance, sequence mapping, technical quality control
        and endpoint-specific mutation labels.
  \item Compare ridge, ordered-context, one-hot and pretrained-representation methods under common
        evaluation conditions; retain historical reference experiments with their original scope.
  \item Test coverage, transfer, head architecture and encoder adaptation using matched controls,
        rather than presuming that greater model complexity improves prediction.
  \item Report regression error, event ranking, PSI-stratum behaviour, seed variation and cluster-level
        uncertainty, with clear distinctions between development and independent validation.
  \item Develop a local React/FastAPI tool for peptide and tiled-protein exploration, source annotation,
        experimental-data comparison, project history and traceable export; verify implementation
        separately from scientific model validity.
\end{enumerate}

\section{Research Subject and Scope}
The research subject is the relation between a single amino acid substitution and a measured change
in peptide-reporter stability. A \emph{residue} denotes one amino acid at a sequence position; a
mutation-effect \emph{footprint} is the map of effects over positions and substitutions, conditional on
the assay and any stated effect threshold.
\begin{itemize}
  \item The main labels come from the human 28-residue peptide study of Zhang et al.\ \citep{zhang2023}.
        The current full-PSI saturation development population contains \nVthreeNPeptides{} peptides
        in \nVthreeNGroups{} protein groups. Historical restricted populations are reported separately;
        the enlarged population is not an independent replication of them.
  \item Sequence, substitution position and amino-acid identity are prediction inputs. Measured WT
        PSI is used to construct labels and define training arms or analysis strata; it is not an input to
        the current sequence-only model. Read count, fold and experimental disagreement are likewise
        kept outside prediction features.
  \item An $L\times20$ display organises predicted substitutions. WT self-substitution is a reference
        entry, not an additional experimentally observed mutation. Missing observations remain missing.
  \item The first software release targets peptide batches, 28-residue whole-protein tiling,
        ClinVar/reference annotation, CSV/TSV comparison, persistent local projects and export.
        These are partly implemented requirements. Scoring tiles across a protein does not establish
        full-length protein effects; imported measurements do not automatically become training data.
  \item The software scope is one local application. Public deployment and accounts are outside the
        current acceptance scope. E3 identification, clinical pathogenicity and therapeutic benefit are
        not model outputs. No new wet-lab experiments have been performed in this project.
\end{itemize}

\section{Research Process and Technologies}
The work proceeds through literature and data audits, hypothesis-driven comparisons, interpretation
of negative as well as positive findings, and prototype development. Python, NumPy, pandas,
scikit-learn and PyTorch support modelling; MMseqs2 supports similarity analysis
\citep{steinegger2017}. Git and \LaTeX{} preserve code and report history. The application design uses
React/TypeScript, FastAPI and local PostgreSQL with a durable inference worker; SQLite supports
isolated baseline development and tests. The implementation status of this design is reported
separately in Chapters~\ref{chap:design} and~\ref{chap:evaluation}.

\section{Contributions Achieved and Still Expected}
As of 7 November 2026, the inspectable contributions are a source-traceable label/QC pipeline,
protein-group-aware development evaluation, controlled comparisons of training coverage and
representations, completed local frozen/LoRA out-of-fold analyses, and a direct equal-budget
retrospective comparison of mutation prioritisation. The report also preserves
negative transfer and weighting findings, and a held-source W5 check that exposes regression and
calibration limits. A restricted-domain baseline and local web/API implementation provide software
evidence within their tested scope.

These are progress contributions in data handling, evaluation and empirical understanding. A
universally optimal architecture, an independent external advantage, a completed expanded software
release and a measured wet-lab benefit remain unestablished. Final model selection and serving
promotion require the additional evidence described in the evaluation and future-work sections.

\section{Structure and Status of This Report}
This is a research progress report built on the existing thesis draft. Chapter~1 specifies the problem
and intended software behaviour; Chapter~2 explains the biology, measurements and computational
concepts; Chapter~3 links research decisions to data and experimental methods; Chapter~4 reports
results, uncertainty and current evidence limits. The concluding part gives provisional findings and
next steps. Appendices retain reproducibility and software instructions. Uncompleted administrative,
experimental and implementation items are not treated as final thesis results.
